% !TEX root = 研究计划.tex
\documentclass[UTF8,zihao=-4,a4paper,fontset=fandol]{ctexart}
\usepackage[margin=24mm,headheight=15pt]{geometry}
\usepackage{amsmath,amssymb,bm,booktabs,array,tabularx,longtable}
\usepackage{graphicx,xcolor,enumitem,fancyhdr,hyperref,listings}
\hypersetup{colorlinks=true,linkcolor=blue!45!black,citecolor=blue!45!black,urlcolor=blue!45!black,breaklinks=true,hypertexnames=false}
\setcounter{tocdepth}{1}
\hypersetup{bookmarksdepth=2}
\setlength{\parindent}{2em}
\setlength{\parskip}{3pt}
\linespread{1.22}
\setlength{\emergencystretch}{2em}
\setlist{nosep,leftmargin=2em,topsep=4pt}
\renewcommand{\arraystretch}{1.18}
\setlength{\tabcolsep}{4pt}
\newcolumntype{Y}{>{\raggedright\arraybackslash}X}
\newcolumntype{P}[1]{>{\raggedright\arraybackslash}p{#1}}
\newcommand{\R}{\mathbb R}
\newcommand{\ind}{\mathbf 1}
\newcommand{\sg}{\operatorname{sg}}
\newcommand{\code}[1]{\texttt{\detokenize{#1}}}
\newcommand{\smalltable}{\small\renewcommand{\arraystretch}{1.18}}
\lstset{basicstyle=\ttfamily\small,breaklines=true,columns=fullflexible,keepspaces=true,showstringspaces=false,frame=single,rulecolor=\color{black!20},xleftmargin=0.5em,xrightmargin=0.5em}
\pagestyle{fancy}\fancyhf{}
\fancyhead[L]{\small ParallelBayes}\fancyhead[R]{\small 中文 v4 一致版}
\fancyfoot[C]{\thepage}
\ctexset{section={format=\Large\bfseries},subsection={format=\large\bfseries},subsubsection={format=\normalsize\bfseries}}
% Generated from a verified scalar bundle; do not edit numbers by hand.
\newcommand{\PBCompactValid}{3741}
\newcommand{\PBCompactNumerical}{42}
\newcommand{\PBCompactUnclassified}{105}
\newcommand{\PBCompactFunctionRows}{14981}
\newcommand{\PBCompactHighRhat}{5000}
\newcommand{\PBCompactUndefinedRhat}{1290}
\newcommand{\PBCPURWMMin}{0.424}
\newcommand{\PBCPURWMMax}{1.178}
\newcommand{\PBCPURWMAbove}{7}
\newcommand{\PBCUDARWMMin}{0.953}
\newcommand{\PBCUDARWMMax}{1.659}
\newcommand{\PBCUDARWMAbove}{14}
\newcommand{\PBCPUMALAMin}{0.116}
\newcommand{\PBCPUMALAMax}{0.394}
\newcommand{\PBCPUMALAAbove}{0}
\newcommand{\PBCUDAMALAMin}{0.234}
\newcommand{\PBCUDAMALAMax}{0.408}
\newcommand{\PBCUDAMALAAbove}{0}


\hypersetup{pdftitle={ParallelBayes 研究与开发计划：中文v4一致版},pdfauthor={ParallelBayes项目}}
\begin{document}
\begin{titlepage}\centering
\vspace*{18mm}
{\LARGE\bfseries ParallelBayes\par}
\vspace{10mm}
{\huge\bfseries 研究与开发计划\par}
\vspace{5mm}
{\Large 时间并行 MCMC 的可核验实现与计算比较\par}
\vspace{12mm}
{\large 中文 v4 一致版\par}
{\large 2026年10月12日\par}
\vspace{16mm}
\begin{minipage}{.87\linewidth}
\textbf{文档用途。}整理当前研究问题、软件范围、数学依据、已执行实验设计与投稿前工作。它是原研究计划的执行后更新，不是重新登记实验，也不是新一轮大规模采样承诺。

\textbf{科学依据。}封存的中文 v4 正文、补充材料及其证据清单。正文18页、补充材料45页、完整结果附录67页保持原样。本版不增加正式重复、GPU测量或新的参考随机点。

\textbf{当前状态。}研究计算与内部技术验收完成；作者信息、材料公开范围、实际审稿访问路线及期刊适配仍需作者确认。当前文档不等于投稿已经获准。
\end{minipage}
\vfill
{\small 固定证据，明确能力，区分执行速度与推断价值。\par}
\end{titlepage}
\pagenumbering{roman}
\tableofcontents\clearpage
\section*{术语与缩写}
\addcontentsline{toc}{section}{术语与缩写}
\noindent\begin{tabularx}{\linewidth}{P{34mm}Y}\toprule
缩写/术语 & 本文含义\\\midrule
MCMC & 马尔可夫链蒙特卡洛\\
RWM & 随机游走 Metropolis 核\\
MALA & Metropolis 校正 Langevin 核\\
NUTS & 无折返采样器；此处为成熟外部基线\\
MSE / RMSE & 均方误差 / 均方根误差\\
MCSE & 蒙特卡洛标准误，描述有限计算的不确定性\\
ESS & 有效样本量；不等于独立实验重复数\\
ESJD & 期望平方跳跃距离；开发选择使用其样本估计\\
JVP & 雅可比—向量积；quasi-DEER 的代理导数计算\\
SBC & 基于模拟的校准，重复先验生成与数据拟合\\
BCa & 偏差校正与加速的 bootstrap 区间\\
IID & 独立同分布；稀有事件例中只作明确假想参照\\
数值有效 & 通过指定数组、路径及输出契约检查，不自动等于后验探索充分\\
参考平方差 & 相对有限数值/采样参考的平方差，不自动等于真值 MSE\\\bottomrule
\end{tabularx}
\clearpage

\pagenumbering{arabic}

\section{研究目标与当前阶段}
\subsection{研究问题}
马尔可夫链蒙特卡洛（MCMC）通常逐步生成下一状态。多链并行把资源分配给多条轨迹；时间并行尝试在固定随机输入下同时求解一段轨迹。后者可以缩短执行深度，却可能增加目标求值、代理导数、扫描与同步开销。更短的运行时间也未必对应更充分的后验探索。

本项目围绕两个可分别检验的问题：第一，固定模型、初值、核参数和实际随机数组时，时间执行器是否降低完成同一递推的成本；第二，在各核采用合理配置后，有限预算内的后验函数误差与完整工作流成本如何变化。前者控制统计信息，后者同时受到初始化、调参和混合影响。数值路径核验、目标分布性质与有限样本精度分开报告，不以其中一项代替另一项。

\subsection{成果定位}
项目形成研究软件、计算比较及可重建材料，不开发替代 Stan/JAGS 的通用概率编程系统，也未提出新的时间并行收敛定理。研究贡献包括可重复使用的目标—核—执行器分离、实际随机输入与独立参考的核对机制，以及由这些机制揭示的重复计算成本和后验探索差异。

现有 Newton 类路径求解及 Picard 方法提供直接技术基础\cite{zoltowski,grazzi,gonzalez}；Stan/BridgeStan、BlackJAX、Pyro 提供模型或基线组件\cite{stan,bridgestan,blackjax,pyro}。compareMCMCs 已提供跨引擎比较，ParallelMCMC.jl 已有时间并行 MALA，BayesForge 已连接 R 与加速推断\cite{comparemcmcs,parallelmcmc,bayesforge}。因此，仅统一接口或调用 GPU 不是本项目的主要创新。

\subsection{范围收束与执行后更新}
最初项目从并行贝叶斯综述转向研究软件与基准。历史 Mac CPU 阶段提供正确性与负结果基础，Windows 原生阶段增加 CPU/CUDA 比较。后来因硬件与磁盘条件收缩正式网格，采用新的紧凑协议身份，同时保留此前已启动或完成的记录。当前正式设计不是在所有计算之前一次性预注册，不能把执行后的计划重写称为前瞻注册。

中文 v4 已完成主研究、有限 H1/L2/W1 补充及 CPU NUTS 定位。下一阶段以整理和投稿为主；自动配置选择器、通用语言翻译、GPU NUTS、融合设备控制及任意 Stan 时间并行均不作为本轮结项前提。现有结果不要求所有压力目标收敛，也不要求出现正加速。

\section{软件方案与支持边界}
\subsection{分层架构}
\begin{table}[htbp]\centering\smalltable
\begin{tabularx}{\linewidth}{P{27mm}Y Y}\toprule
层次 & 职责 & 必须记录的信息\\\midrule
目标与坐标 & 对数密度、导数、参数变换与输出函数 & 数据、坐标、参数顺序、实现哈希\\
转移核 & RWM、MALA及成熟 NUTS 基线 & 步长、初值、预热、真实随机输入\\
执行器 & 顺序、Online Picard、窗口 quasi-DEER & 窗口、求解轮数、确认前缀、停止原因\\
资源与计量 & CPU/CUDA、同步、内存及计时 & 线程、依赖、设备、成本边界\\
结果与证据 & 样本、诊断、失败及重建 & 数组、审计、原始状态、完整性清单\\\bottomrule
\end{tabularx}
\caption{模块职责。模型、核和执行方式独立描述，避免以单个“并行 MCMC”参数混淆统计与计算机制。}
\end{table}

R 管理目标、实验调用及 posterior 样本对象；编译的 CPU 模型和 Python/JAX/torch 内核承担数值计算。BridgeStan 通过编译的 C++ 密度接口运行 CPU 顺序核，JAX/torch 只接入明确实现并核验的原生目标。R–Python 边界尽量使用批量调用，其导入、传输与转换成本必须保留。
\begin{table}[htbp]\centering\smalltable
\begin{tabularx}{\linewidth}{P{36mm}Y Y Y}\toprule
目标入口及设备 & 顺序 RWM/MALA & Picard RWM、quasi-DEER MALA & NUTS\\\midrule
Stan/BridgeStan，Mac CPU & 支持 & 不支持 & 不支持\\
明确原生 JAX 目标，CPU & 支持 & 支持对应组合 & BlackJAX 基线\\
明确原生 torch 目标，Windows CPU & 已核验 & 已核验，含主机控制 & 独立 Pyro 研究 CLI\\
明确原生 torch 目标，Windows CUDA & 已核验 & 已核验，含主机控制 & 未接入或核验\\\bottomrule
\end{tabularx}
\caption{v4 对应的软件能力与实际设备边界。原生 Windows Stan 编译尚未核验。Pyro 四进程基线不属于通用 \code{pb_sample()} 的 torch NUTS 选项，也不是 NUTS 轨迹的时间并行。}
\end{table}

\code{pb_capabilities()} 查询包接口组合；独立研究 CLI 的契约另列。\code{pb_benchmark()} 运行配置列表并组织 R 结果，完整冻结协议、任务恢复和原始证据管理由 Python CLI 及项目文件完成。Stan 入口没有默认独立轨迹参考，\code{audit=TRUE} 明确拒绝。扩展目标需提供原生 Python 模型及独立参考，不支持把任意 Stan 文件自动编译为 JAX/torch GPU 目标，也不是任意 R 闭包的一键时间并行接口。


\subsection{最小公开使用流程}
真实运行的包接口示例如下；正式 CLI 网格另由协议驱动。
\begin{lstlisting}[language=R]
library(parallelbayes)
model <- pb_model("gaussian", dimension = 2L, backend = "torch")
stopifnot(pb_validate(model)$passed)
pb_capabilities(model)
configs <- lapply(c("sequential", "online_picard"), function(executor)
  list(kernel = "rwm", executor = executor, device = "cpu",
       draws = 128L, chains = 4L, step_size = 0.4,
       window = 16L, seed = 2026L, audit = TRUE))
fits <- pb_benchmark(model, configs)
\end{lstlisting}
同核两次调用产生相同实际输入与接受数组，并通过独立参考。短例只展示接口与核验，不能凭128步宣称后验收敛或加速。新增模型的示例采用 Poisson 目标及独立 NumPy 参考，由 R 经 reticulate 访问；模型核对先于性能比较。

\section{转移规则与数值核验依据}
\subsection{目标、坐标与实际随机输入}
设原参数为 $\theta=T(q)$，$T$ 是相应自由坐标之间的可微一一变换，实际计算使用 $q\in\R^d$。计算目标为
\begin{equation}\label{eq:target}
\ell(q)=\log p(y,T(q))+\log|\det J_T(q)|.
\end{equation}
约束变换、参数顺序和 Jacobian 项共同决定目标分布。两个实现若省略不同的参数无关常数，核对 $\ell(q)-\ell(q_0)$，而非要求原始对数密度相等；梯度必须在同一坐标下比较。正值测试使用 $\theta=e^q$ 及 lognormal 密度，并通过故意遗漏 Jacobian 检查测试系统的敏感性。

固定初值、核参数、提议噪声和均匀随机数组后，MCMC 成为确定递推 $q_t^\star=F_t(q_{t-1}^\star)$。配对执行器读取同一实际数组；两个软件设置同一整数种子并不能替代这项控制。采样数组由锁定 NumPy 版本的 Philox 生成，导数探针另用随机流，数组及哈希一并保存。独立 NumPy 顺序参考使用独立密度、解析梯度和显式循环，不调用待核验的 JAX 或 torch 转移函数。

\subsection{两种转移核}
随机游走 Metropolis（RWM）使用 $q'=q+s z$，$z\sim N(0,I_d)$，对数接受比为 $a=\ell(q')-\ell(q)$。Metropolis 校正 Langevin 算法（MALA）使用
\begin{align}
q'&=q+h\nabla\ell(q)+\sqrt{2h}\,z,\\
a(q,q')&=\ell(q')-\ell(q)-
\frac{\|q-q'-h\nabla\ell(q')\|^2-\|q'-q-h\nabla\ell(q)\|^2}{4h}.
\end{align}
当 $\log u<a$ 时接受，否则保留 $q$；拒绝自环仍计为一次转移。这里 $h$ 不是提议噪声的标准差。NUTS 基线保留自身随机流与适应预热，不与 MH 核共享整条路径。

\subsection{窗口 quasi-DEER MALA}
在固定输入下，窗口内第 $k+1$ 轮路径满足
\begin{equation}\label{eq:deer}
q_t^{(k+1)}=F_t(q_{t-1}^{(k)})+
B_t^{(k)}\{q_{t-1}^{(k+1)}-q_{t-1}^{(k)}\}.
\end{equation}
旧路径上的映射和 $B_t^{(k)}$ 可批量计算，再将右侧整理为仿射递推，通过结合扫描求解。精确算术下，若各映射有定义、$B$ 有限且扫描被精确求解，第 $k$ 轮至少恢复前 $k$ 个状态，长度 $W$ 的窗口至多经过 $W$ 轮恢复原递推\cite{gonzalez,zoltowski}。固定窗口起点使第一步正确；前 $k$ 步正确时，下一轮保持这段前缀，第 $k+1$ 步的修正项为零。这一因果性质不要求 $B$ 是真实 Jacobian，也不保证浮点实现提前停止后的路径精度。

实际实现将硬接受变量 $b=\ind\{\log u<a\}$ 与软门控组合：
\begin{equation}
g=b+\sigma(a-\log u)-\sg\{\sigma(a-\log u)\},\qquad
B_t=\operatorname{Diag}\!\left[\operatorname{clip}\{v_t\odot(J_{\rm proxy,t}v_t),-1,1\}\right].
\end{equation}
前向仍使用硬 MH 输出；torch 另以停止梯度构造避免软项的舍入抵消。$J_{\rm proxy}$ 是代理 Jacobian，而非硬 MH 映射的真实导数。每链每位置一个 Rademacher 探针 $v_t$，在窗口迭代中复用；截断仅作用于导数近似。

原论文伪代码的软门控使用 $\sigma(\min(0,a)-\log u)$；本项目归档上游提交 \code{5e5b637} 及适配实现使用未截断的 $a$。硬接受判据等价，代理导数一般不同。此外，本研究使用非重叠窗口、重复窗口起点的初猜和硬残差停止，不使用 NaN 归零或状态裁剪。因此性能结论属于所测适配实现，不能视为原实现的直接复现。

\subsection{Online Picard RWM}
对于已经确认的窗口起点 $q_L$，先计算固定噪声对应的接受增量，再通过前缀和重建候选路径\cite{grazzi}：
\begin{align}
D_t^{(k)}&=s z_t\ind\{\log u_t<\ell(q_{t-1}^{(k)}+s z_t)-\ell(q_{t-1}^{(k)})\},\\
q_t^{(k+1)}&=q_L+\sum_{i=L+1}^{t}D_i^{(k)}.
\end{align}
在新路径的前驱处重新计算接受事件，只提交连续一致的最长前缀。若四个位置依次为“一致、一致、不一致、一致”，只能提交前两个位置，第四个位置的前驱尚未确认。torch 多链按各链确认长度的最小值共同推进，末尾窗口使用实际剩余长度。浮点前缀和与顺序加法可能不同，事件匹配仍需配合独立路径核对。

\subsection{数值输出标准}
quasi-DEER 的窗口停止检查硬递推残差；Picard 使用前述连续前缀确认。残差条件为
\begin{equation}
\max_{t,j}\frac{|q_{t,j}-F_t(q_{t-1})_j|}
{\mathrm{atol}+\mathrm{rtol}|F_t(q_{t-1})_j|}\leq1.
\end{equation}
正式审计另按每链检查最大路径差不超过
$100[\mathrm{atol}+\mathrm{rtol}\max(1,\max_{t,j}|q_{t,j}|)]$，其中两项容差均为 $10^{-10}$；接受事件失配必须为零。局部残差小不推出长路径误差小，代理导数也不是后验偏差证明。历史判据使用链内整体状态尺度，后续逐函数分析补充其含义，未追溯修改资格。

非有限值、迭代耗尽和外部路径失配都保留为失败，失败数组不能转为正式后验样本。本版将非有限提议密度统一视为失败，尚未一般区分合法支持集外 $-\infty$ 与 NaN/计算异常；这一处理限制支持范围。主研究未用重抽随机数、删除困难状态或顺序回退改变资格。


\section{已执行的正式研究设计}
\subsection{正式模型库与坐标}
正式研究包括九个目标、36 个预先指定的后验函数。为便于数值核对与解释，区分基础模型坐标 $r$、实际采样坐标 $q$ 和输出参数 $\theta$：$r=c+Lq$，$\theta=T_0(r)$；目标含 $\log|\det L|$ 及所需约束项。

\begin{table}[htbp]\centering\smalltable
\begin{tabularx}{\linewidth}{l r Y Y}\toprule
目标 & 维数 & 模型与数据 & 坐标及比较作用\\\midrule
G1 & 8 & 标准正态 & $c=0,L=I$；解析控制\\
G2 & 64 & 旋转高斯，原协方差条件数100 & 完整 Cholesky 白化；高维稠密计算\\
A1 & 64 & 潜在 AR(1)，$\rho=0.9$，单位观测噪声 & 后验均值平移与完整白化\\
L1 & 8 & logistic，$n=256$，先验标准差2.5 & 众数及观测信息坐标\\
L2 & 8 & logistic，$n=4096$，同一生成系数 & 似然成本及极稀有事件\\
H1 & 9 & 中心化漏斗，$v\sim N(0,9)$ & 原坐标，检验几何与路径敏感性\\
H2 & 9 & 非中心化漏斗 & 输出回原尺度，核随坐标改变\\
M1 & 8 & $\frac12N(-5e_1,I)+\frac12N(5e_1,I)$ & 双峰探索与分散初始化\\
W1 & 2 & 水井真实数据，3020条 & 截距、距离系数；二维参考\\\bottomrule
\end{tabularx}
\caption{当前 v4 的正式目标库。完整数据、矩阵、浮点配置及哈希随冻结协议保存。}
\end{table}

G2 的协方差为 $\Sigma=U\Lambda U^\mathsf T$，$\Lambda_{jj}=100^{(j-1)/63}$，$U$ 来自固定正态矩阵的 QR 分解。A1 的先验协方差为 $K_{ij}=0.9^{|i-j|}$，固定数据后 $\Sigma=(K^{-1}+I)^{-1}$、$m=\Sigma y$。二者在实际 $q$ 坐标均为标准正态，原尺度相关性不能直接解释本轮核的混合差异。

L1/L2 的设计含截距及七列标准化正态变量，系数先验为 $N(0,2.5^2)$，生成系数相同但数据不同。W1 为
$\Pr(y_i=1)=\sigma(\alpha+\beta d_i/100)$，在两个系数上用平坦先验。三种回归目标从零点开始，以 trust-exact 求众数；目标有限且解析梯度最大绝对值不超过 $10^{-7}$ 时接受，最多300步，$L$ 为观测信息逆矩阵的 Cholesky 因子。未做特征值修复，也未用参考样本构造坐标；该变换不把 Laplace 近似当作后验真值。

H1 为 $x_i\mid v\sim N(0,e^v)$，$i=1,\ldots,8$；H2 使用独立 $z_i\sim N(0,1)$ 并输出 $x_i=z_i e^{v/2}$。两者表示同一原尺度后验，但在各自坐标使用同一数值步长并不意味着同一原尺度转移核。

\subsection{独立开发选择与冻结配置}
步长在正式输入之外选择。RWM 候选为 $\{0.5,1,1.5,2.4,3.5,5\}/\sqrt d$；MALA 候选为 $\{0.1,0.3,0.6,1,1.6,2.4\}/d^{1/3}$。每候选两份四链输入，丢弃512步、保留1024步；仅在两次均有效时，以平均 ESJD/$d$ 最大者入选，平局取较小步长。该准则衡量局部移动，未优化 ESS 或稀有事件精度。

\begin{table}[htbp]\centering\smalltable
\begin{tabular}{lrrl}\toprule
目标 & MALA $h$ & RWM $s$ & 固定中心及线性变换\\\midrule
G1 & 0.8 & 0.848528137 & $0,I$\\
G2 & 0.4 & 0.3 & $0,\operatorname{chol}(\Sigma)$\\
A1 & 0.4 & 0.3 & $m,\operatorname{chol}(\Sigma)$\\
L1/L2 & 0.8 & 0.848528137 & 众数及观测信息\\
H1 & 0.480749857 & 1.66666667 & $0,I$\\
H2 & 0.769199771 & 0.8 & $0,I$\\
M1 & 0.8 & 0.848528137 & $0,I$\\
W1 & 1.26992084 & 1.69705627 & 众数及观测信息\\\bottomrule
\end{tabular}
\caption{冻结配置的显示精度。归档 JSON 提供完整精度；$h$ 对应噪声 $\sqrt{2h}$，$s$ 是 RWM 提议标准差。}
\end{table}

开发批次216项中212项完成，4项 H1 MALA 失败保留。H1 RWM 选择了候选上边界，说明配置仍受有限搜索范围约束。已记录开发任务内部成本226.917秒、坐标构造0.128秒；其中包含审计及写盘，未包含新进程导入，不能直接并入或扣除正式普通调用成本。

\subsection{正式网格、预算与资源}
每目标24份独立四链实际输入；同核的两个执行器及两设备读取同一输入。九种工作流为 CPU/CUDA 的四种 MH 组合，以及独立 CPU Pyro NUTS。两档预算每链保留1024或4096次转移，总计
\begin{equation}9\times9\times2\times24=3888\ \text{项主任务}.
\end{equation}
四链、时间执行器和设备重复不增加独立重复数。缓存探测限于 G1/G2/L2/W1，共256项、1024次调用；不计入正式后验拟合。

四链初值来自 $N(0,4I)$，M1 首坐标设为 $(-5,5,-5,5)$。MH 与 NUTS 共用初值，但 NUTS 使用独立随机流及预热。MH 丢弃512步、窗口32；quasi-DEER 每窗口最多2048轮，Picard最多 $N+512$ 轮。NUTS 用四个原生 spawn 子进程，每链预热1024步、目标接受率0.8、最大树深8，适应步长及对角质量矩阵，未使用 JIT。

正式 Windows 环境为原生 Windows 11、AMD Ryzen 7 9800X3D、RTX 5080，float64；Python 3.12.14、NumPy 2.2.6、SciPy 1.15.3、torch 2.13.0+cu130、Pyro 1.9.2、psutil 7.1.0，R 4.6.1、posterior 1.7.0。MH 进程 torch 线程4、interop线程1；NUTS每个子进程各为1。BLAS线程环境变量未统一设置，不能把这些数字解释为所有库的硬上限或独占物理核心预算。

存储/工作区估计上限2048 MiB、进程树 RSS 监督上限8 GiB、Windows Job提交内存上限12 GiB，三者计量不同。启动要求可用主存至少12 GiB、磁盘4 GiB，CUDA另需空闲显存3 GiB，运行磁盘底线1 GiB。资源不足暂停启动，数值失败、进程失败与未启动分别记录。没有总实验时限，也未通过重试或回退改变样本资格。

\subsection{函数、尺度与参考}
\begin{table}[htbp]\centering\smalltable
\begin{tabularx}{\linewidth}{l Y Y}\toprule
目标 & 原尺度上的函数集合 & 参考及其性质\\\midrule
G1/G2/A1 & 标准化首参数 $z$、$z^2$、$\ind(z>1)$ & 解析：$0,1,1-\Phi(1)$\\
L1/L2 & $\theta_1,\theta_2,\sigma(\theta_1),\ind(\theta_1>0)$ & 独立有限 NUTS；L2事件另作重要性参考\\
H1/H2 & $v/3,\tanh(v/3),\ind(x_1>0),\cos(x_1e^{-v/2})$ & 解析：$0,0,1/2,e^{-1/2}$\\
M1 & $\theta_1/\sqrt{26},\ind(\theta_1>0),\theta_2,\theta_2^2$ & 解析：$0,1/2,0,1$\\
W1 & $\alpha,\beta,\alpha^2,\beta^2,p(0),p(100),\ind(\beta>0)$ & 冻结求积点；v4另给确定数值包络\\\bottomrule
\end{tabularx}
\caption{全部36函数。各尺度 $s_f=1$，因此阈值表示函数原单位上的绝对误差，而非自动统一的相对精度。}
\end{table}

对方法 $a$ 和预算 $N$，在数值有效且函数有限的重复集合 $V$ 上计算
\begin{equation}\label{eq:loss}
\widehat E_{a,f}(N)=\frac{1}{|V|}\sum_{r\in V}
\frac{(\widehat\mu_{a,f}^{(r)}(N)-\mu_f^{\rm ref})^2}{s_f^2}.
\end{equation}
解析参考时可解释为有效输出条件下的 MSE 估计；有限参考时称为参考平方差。比较两方法固定共同有效集合，同时报告全部24份的失败分母，不以 $\hat R$ 达标筛掉困难拟合。预算点上的最大函数平方差 $\leq0.01$ 对应 RMSE $\leq0.1$，只对指定函数与尺度有意义。误差估计的 MCSE 判据与这一误差阈值分开，不把短链 ESS 当作真实 MSE。

主分析使用9999次逐点 BCa bootstrap，单位为整份四链独立输入，至少20份共同有效重复才计算区间；退化情况保留未定，不作全体比较的同时置信声明。秩正态化与折叠 $\hat R$、bulk/tail ESS 由 posterior 计算\cite{vehtari,posterior}；常量函数不填成 $\hat R=1$。预算仅两个点，报告已测预算的误差与耗时，不插值为精确达标时间。

L1/L2 每模型四份独立四链 NUTS 参考，每链预热1024、保留16384步，共262144个相关状态；MCSE取批均值与拟合间估计的较大者。无记录发散不能排除共同偏差。L2原符号事件全零，其不确定性保持未定。W1 原参考以12个保存的浮点求积配置检查稳定性；这些变化不是认证误差界。v4的 L2 重要性参考和 W1 确定包络单独报告，不覆盖原参考身份。

\subsection{成本边界与比较单位}
\begin{table}[htbp]\centering\smalltable
\begin{tabularx}{\linewidth}{l Y Y}\toprule
口径 & 包含内容 & 回答的问题\\\midrule
缓存执行 & 已准备内核的执行，显式设备同步 & 同一指定计算能否更快完成？\\
普通批量工作流 & 进程启动、输入/目标准备、预热/丢弃、采样、变换、正常结果和 R 诊断 & 获得可分析输出的实际批量成本\\
研究审计 & 再含独立路径参考、保护、完整证据归档 & 取得可核验研究材料的成本\\
R 前端技术成本 & 最外层 R 启动/依赖、结果传回与转换的实测 & R 使用者首次与重复调用的边界\\\bottomrule
\end{tabularx}
\caption{成本口径相互嵌套或属于另测边界，不能直接相加。开发调参、整库环境安装和全局归档上传另列。}
\end{table}

同设备、同核、同输入配对的速度比为 $S=T_{\rm seq}/T_{\rm par}$，$S>1$ 表示时间执行更快。不同核需结合函数误差与诊断解释。跨 Mac/JAX 和 Windows/torch 的差异同时包含硬件、后端及实现，不归因于算法单一因素。批处理向量化不等于已证明物理核心有效利用；CPU/CUDA计时包含实际同步及主机控制成本。


\section{验证路线、实验记录与判断规则}
\subsection{正确性先于性能}
验证分为四层。模型层核对同坐标密度差、梯度、变换及 Jacobian；执行层核对独立顺序参考、状态路径及实际接受数组；统计层使用解析函数和基于模拟的校准（SBC）；实验层记录全部失败与资源状态。有限测试不能替代一般数学证明，但可以定位实现错误与比较口径混淆。

SBC 采用 $\theta\sim N(0,1)$、$y_j\mid\theta\sim N(\theta,1)$、$j=1,\ldots,8$，解析后验为 $N(\sum_jy_j/9,1/9)$。同一批64组真参数—数据用于各工作流，并与解析区间逐数据集配对，不能把五条工作流当作五份独立校准证据。SBC 测试量的敏感性也取决于所选函数\cite{modrak}。

错误注入与退化输入用于检查遗漏 Jacobian、接受分支错误、常量函数及非有限失败能否被记录。恢复任务必须匹配源码、协议、数据与校验和；匹配时复用完成输出，身份不一致时停止，而非覆盖冻结结果。

\subsection{参考精度与补充协议}
H1 补充只在原路径上比较函数差，不能用小函数均值差重新认可失败轨迹。L2 用两套冻结重要性提议及独立批次评估极稀有符号事件，MCSE 是有限近似而非认证界\cite{owen}。W1 使用任意精度球算术、两条有限域求积路线与显式尾部控制建立数值包络\cite{arb,dlmfquadrature}，确定数值误差与重复采样误差分别报告。CPU NUTS 定位通过预先安排的诊断开关和 worker 数检查资源路径，不追溯解释全部历史故障。

以上新增分析与参考均有独立身份，不替换原正式参考或主任务资格。原全零事件、未分类失败、失败部分链及负结果继续保留。均值接近、诊断良好、接受事件一致和设备探测成功分别回答不同问题。

\subsection{解释规则与适用边界}
同核配对首先解释执行成本，跨核比较再解释固定预算的函数误差。报告全部失败的分母、条件可用集合及参考来源；区间基于真实重复数，不把更多 bootstrap 重抽或缓存调用视为增加实验信息。常量稀有事件的极小绝对误差与100\%相对误差可以同时成立。多峰目标均衡初值产生的良好池化均值也可能掩盖缺少跨峰探索。

自动配置选择器只有在未见模型族上、计入探测与额外编译后优于合理固定配置和简单搜索，才进入主要贡献。当前未开展这项正式验证，也不为本轮论文补造自动推荐结果。

\section{工作包完成情况与剩余安排}
\begin{longtable}{P{24mm}P{57mm}P{39mm}P{25mm}}
\caption{研究计划的实际执行状态。完成指当前规定范围内的产物与证据齐备。}\\
\toprule 工作包 & 完成产物及核验 & 保留的范围限制 & 状态\\\midrule\endfirsthead
\toprule 工作包 & 完成产物及核验 & 保留的范围限制 & 状态\\\midrule\endhead
目标与顺序参考 & 高斯、logistic、约束测试；独立 NumPy；实际数组管理 & 非通用模型编译器 & 已完成\\
两类执行器 & 窗口 quasi-DEER、Online Picard；残差/分支/前缀与失败记录 & 针对适配实现；有路径失败 & 已完成\\
R接口与基线 & 能力查询、采样、验证、基准入口；JAX与原生torch；CPU NUTS CLI & Stan只CPU顺序；GPU NUTS未核验 & 已完成\\
正式研究 & 3888主任务、256缓存；独立接收、4144分析与恢复检查 & 两个预算；105未分类失败保留 & 已完成\\
有限补充 & H1函数路径、L2重要性参考、W1数值包络、S2资源定位 & 不追溯改变主实验 & 已完成\\
R普通调用 & 原32进程64次技术调用；安装新目录示例 & 技术计时非新正式重复 & 已完成\\
复现封存 & 41组件清单、依赖闭包、源码与表图重建、v4 PDF & 同一Mac内部新环境；非外部团队 & 已完成内部验收\\
投稿适配 & 作者声明、访问路线、目标期刊及后继版本 & 需要真实作者决定 & 待确认后实施\\\bottomrule
\end{longtable}

当前不再启动“为了出现加速”而寻找案例的新网格。若作者日后决定拓展到昂贵似然、融合执行或配置选择，应另设问题、预算和协议，并单独报告结果。本计划没有隐含的新增必做实验，也不重新运行已经通过且未受修改影响的采样检查。

后续按材料、访问、模板、终审顺序推进：先从封存候选建立后继投稿目录，再根据作者授权安排真实材料获取测试；随后按目标期刊调整语言、篇幅和声明；最后由作者确认投稿版本。作者或期刊提出实质方法变更时，评估其是否需要新的验证身份，避免与原研究混写。

\section{复现材料、版本及交付管理}
\subsection{源码与证据分别管理}
Git 跟踪第一方源码、测试、协议、小型汇总、论文及图件；大原始数组单独归档，以精确成员、大小和 SHA256 清单连接。第三方论文/书籍全文、私有技能、环境缓存及凭据不属于 Git 资产。原 Mac 0.1.0/0.1.1、冻结协议和历史结果保存，不用当前源码假代旧实验性能。

紧凑正式执行源为 \code{3a37a89}，交付源 \code{94bebec}；中文 v4 论文构建源 \code{699e645}，投稿候选 \code{b84e966}、标签 \code{submission-v4-rc1}。不同身份用于不同证据，不相互替代。当前文档在后继文档分支新增，不修改这些冻结对象。

\subsection{内部验收与可用性}
统一入口由 \code{docs/SUBMISSION-REPRODUCE.md} 和论文—证据—命令映射组织。41组件列示39,354,532,214字节，约36.65 GiB；历史包存在重叠，不代表独立信息量。轻量入口包包含源码、论文输入、安装候选、清单及验收记录，不包含全部原始大组件。

新分析环境从4144份保存分析记录重算统计，4589个 bootstrap 数组一致；从新解包目录重建81个主 CSV、4个补充 CSV、6个主文片段与45幅图。169个 LaTeX 内容输入、正文与 SI 全63页的文本及72 dpi像素保持一致。新软件安装核验17个模块并完成两次 CPU 技术调用，未重测 GPU。原来的全量路径重放证据继续提供，本轮保存记录重聚合不冒充再次全量重放。

内部验收曾暴露依赖锁遗漏，首次失败日志保留；补齐真实依赖而不修改冻结采样核心。这些证据来自同一 Mac 的新环境与目录，R依赖库明确复用原锁定库；不能声称是异机或非作者团队复现。

\section{需作者确认的投稿事项}
\begin{enumerate}
\item 作者姓名、顺序、单位、通讯邮箱、ORCID及贡献分工。不能从 Git 身份推断学术署名。
\item 资助项目与编号、利益冲突、致谢对象；无真实信息时保留待填，不虚构声明。
\item 哪些证据可公开、哪些只供审稿，以及本地 H1/L2/W1 原件的分享路线。此前允许追加私有草稿不等于允许公开全部材料。
\item 第一方数组及第三方水井数据的权利、许可和机构要求。代码许可不自动扩展为全部数据许可。
\item 审稿人实际可取得材料的固定入口或受控流程。草稿 Release 地址本身不是匿名访问验收。
\item 目标期刊、文章类型、语言、匿名要求与 AI 使用披露。确认后按官方规则适配，再封存后继版本。
\item 作者对最终科学结论、失败与不确定性边界及材料清单的终审批准。项目不自动向期刊或 CRAN 提交。
\end{enumerate}

这些事项影响投稿，而不意味着主研究计算仍未结束。结项标准是软件能力清楚、指定输出有据可查、负结果和失败可重建、结论范围准确；不是证明所有算法收敛或 CPU/GPU 普遍加速。

\clearpage% Selected unchanged entries from frozen v4; scientific usage is described in the documents.
\begin{thebibliography}{99}\small\setlength{\itemsep}{3pt}
\bibitem{zoltowski} Zoltowski D.M., Wu S., Gonzalez X., et al. \href{https://proceedings.nips.cc/paper_files/paper/2025/hash/202886ee1c9ca735cb5bff3a00a69883-Abstract-Conference.html}{Parallelizing MCMC Across the Sequence Length}. Advances in Neural Information Processing Systems, 2025, 38.
\bibitem{grazzi} Grazzi S., Zanella G. \href{https://doi.org/10.1093/biomet/asag022}{Parallel computations for Metropolis Markov chains with Picard maps}. Biometrika, 2026, 113(2): asag022.
\bibitem{stan} Carpenter B., Gelman A., Hoffman M.D., et al. \href{https://doi.org/10.18637/jss.v076.i01}{Stan: A Probabilistic Programming Language}. Journal of Statistical Software, 2017, 76(1): 1--32.
\bibitem{bridgestan} Roualdes E.A., Ward B., Carpenter B., Seyboldt A., Axen S.D. \href{https://doi.org/10.21105/joss.05236}{BridgeStan: Efficient in-memory access to the methods of a Stan model}. Journal of Open Source Software, 2023, 8(87):5236.
\bibitem{blackjax} Cabezas A., Corenflos A., Lao J., Louf R., et al. \href{https://arxiv.org/abs/2402.10797}{BlackJAX: Composable Bayesian inference in JAX}. arXiv:2402.10797, 2024.
\bibitem{pyro} Bingham E., Chen J.P., Jankowiak M., et al. \href{https://jmlr.org/papers/v20/18-403.html}{Pyro: Deep Universal Probabilistic Programming}. Journal of Machine Learning Research, 2019, 20(28):1--6.
\bibitem{comparemcmcs} NIMBLE development team. \href{https://github.com/nimble-dev/compareMCMCs}{compareMCMCs: Tools for comparing MCMC efficiency from NIMBLE and other engines}. 软件作者仓库，访问于2026年10月3日。
\bibitem{posterior} Stan development team. \href{https://mc-stan.org/posterior/reference/rhat.html}{posterior: Rhat, bulk ESS and tail ESS documentation}. 软件文档，版本1.7.0。
\bibitem{parallelmcmc} Senne R. and contributors. \href{https://github.com/rsenne/ParallelMCMC.jl/tree/7be9c866c8be03ad994a5e7f0486ec3290088813}{ParallelMCMC.jl}. 软件仓库，提交7be9c86，访问于2026年10月6日。
\bibitem{bayesforge} Sosa S. \href{https://cran.r-universe.dev/BayesForge/doc/manual.html}{BayesForge: Bayesian Inference using numpyro and XLA}. R包0.0.1参考手册，访问于2026年10月6日。
\bibitem{gonzalez} Gonzalez X., Warrington A., Smith J.T.H., Linderman S.W. \href{https://arxiv.org/abs/2407.19115v3}{Towards Scalable and Stable Parallelization of Nonlinear RNNs}. arXiv:2407.19115v3, 2025（初版2024）.
\bibitem{vehtari} Vehtari A., Gelman A., Simpson D., Carpenter B., Bürkner P.-C. \href{https://doi.org/10.1214/20-BA1221}{Rank-Normalization, Folding, and Localization: An Improved R-hat for Assessing Convergence of MCMC}. Bayesian Analysis, 2021, 16(2):667--718.
\bibitem{modrak} Modrák M., Moon A.H., Kim S., et al. \href{https://doi.org/10.1214/23-BA1404}{Simulation-Based Calibration Checking for Bayesian Computation: The Choice of Test Quantities Shapes Sensitivity}. Bayesian Analysis, 2025, 20(2):461--488. DOI:10.1214/23-BA1404.
\bibitem{owen} Owen A.B. \href{https://artowen.su.domains/mc/}{Monte Carlo theory, methods and examples}. 2013, Chapter 9: Importance sampling. 在线书稿。
\bibitem{dlmfquadrature} National Institute of Standards and Technology. \href{https://dlmf.nist.gov/3.5}{Digital Library of Mathematical Functions, Section 3.5: Quadrature}. Equations 3.5.6, 3.5.19--3.5.21，访问于2026年10月10日。
\bibitem{arb} Johansson F. \href{https://doi.org/10.1109/TC.2017.2690633}{Arb: Efficient Arbitrary-Precision Midpoint-Radius Interval Arithmetic}. IEEE Transactions on Computers, 2017, 66(8):1281--1292.
\end{thebibliography}

\end{document}
